% =====================================================================
% ANALISIS DE IMPACTO -- Responsable: Fabian Moreno Ugarte
% Orden solicitado: etico, social, economico, ambiental, sintesis.
% =====================================================================
\section{Análisis de impacto}
\label{sec:analisis}

Esta sección evalúa el impacto del sistema en cuatro dimensiones ---
ética, social, económica y ambiental --- y cierra con una síntesis que
las integra. Su estructura --- identificar el tratamiento, valorar
necesidad y proporcionalidad, identificar riesgos para los derechos de
las personas y proponer medidas de mitigación --- reproduce
deliberadamente la de una Evaluación de Impacto en Protección de Datos.

Esto no es una analogía metodológica. Como se expone en la
Sección~\ref{sec:eipd}, el D.S. 016-2024-JUS exige una EIPD antes de
iniciar tratamientos de alto riesgo, y la videovigilancia masiva figura
expresamente entre los supuestos que la activan. Esta sección es, por
tanto, el insumo sustantivo de un documento que el proyecto está obligado
a producir, y no un ejercicio analítico de acompañamiento.

Un criterio guía la redacción: se prefiere declarar un impacto como no
cuantificado antes que estimarlo sin base. Varias de las magnitudes que
naturalmente irían en esta sección --- número de pasajeros afectados,
costo evitado, consumo energético --- dependen de información operativa
de LAP o de mediciones que el equipo aún no ha realizado. Donde ocurre,
se indica expresamente.

% =====================================================================
\subsection{Impacto ético}
\label{sec:impacto-etico}

\subsubsection{Asimetría entre el sujeto observado y el operador}

El pasajero de un aeropuerto no elige ser observado. No puede negociar
las condiciones, no puede optar por un terminal sin cámaras y, en la
práctica, tampoco puede optar por no volar. Esta ausencia de alternativa
razonable es lo que hace que el consentimiento sea una base inadecuada
para este tratamiento: un consentimiento que no puede negarse sin
renunciar al servicio no es libre.

De ahí se sigue una consecuencia que conviene enunciar con claridad: si
el consentimiento no es la base habilitante, la legitimidad del sistema
descansa enteramente en que supere el examen de proporcionalidad y en que
el deber de información se cumpla de manera efectiva. El cartel del
Anexo~\ref{anx:cartel} no es un trámite: es, junto con la
proporcionalidad, uno de los dos pilares sobre los que se sostiene la
licitud del tratamiento.

\subsubsection{Riesgo de deslizamiento de finalidad}

Este es, a juicio de este frente, el riesgo ético principal del proyecto,
y no proviene de un defecto técnico sino del éxito del sistema.

Una vez desplegada la infraestructura --- cámaras calibradas, canal de
video, capacidad de inferencia, tablero de visualización --- el costo
marginal de añadir capacidades es bajo. Añadir reidentificación para
medir tiempos de permanencia, o reconocimiento facial para detectar
personas de interés, es una decisión de configuración y de presupuesto,
no un rediseño. La barrera que impide ese deslizamiento no es técnica: es
documental y contractual.

\cajaaviso{Medida de mitigación propuesta}{Documentar la finalidad de
forma restrictiva y explícita, e incorporar en el entregable una
declaración de finalidades excluidas: el sistema no realiza
identificación de personas, no realiza seguimiento individual de
trayectorias, no realiza reconocimiento de emociones y no alimenta
decisiones automatizadas sobre pasajeros individuales. Una lista de lo
que el sistema \emph{no} hace es más difícil de erosionar silenciosamente
que una lista de lo que sí hace.

Este frente propone que esa declaración se incorpore como entregable
formal del proyecto. Es una de las contribuciones más concretas que puede
aportar, y su valor depende precisamente de que sea explícita: una
finalidad excluida que solo consta en el ánimo del equipo no restringe
nada.}

\subsubsection{Sesgo y desempeño desigual}

Ninguno de los tres artículos revisados reporta desempeño desagregado por
características demográficas de las personas contadas. No se trata de una
omisión de este informe sino de una ausencia en la literatura consultada,
y tiene consecuencias concretas para un aeropuerto internacional, donde la
población observada es considerablemente más heterogénea que la de los
conjuntos de entrenamiento disponibles.

Los factores que razonablemente podrían inducir desempeño desigual son
identificables aunque no estén cuantificados:

\begin{itemize}[leftmargin=1.5em,itemsep=0.3em]
  \item \textbf{Estatura.} Un contador entrenado con anotaciones de
        cabezas puede subcontar sistemáticamente a niños, especialmente
        en escenas densas donde quedan ocluidos por adultos.
  \item \textbf{Vestimenta y cubrimientos de cabeza.} Sombreros, gorras,
        capuchas y prendas religiosas alteran la apariencia de la región
        de la cabeza, que es precisamente la señal que estos métodos
        explotan.
  \item \textbf{Equipaje y sillas de ruedas.} Un pasajero con equipaje
        voluminoso o en silla de ruedas presenta una silueta distinta de
        la típica en los conjuntos de entrenamiento.
  \item \textbf{Composición demográfica del conjunto de origen.} Los
        conjuntos de referencia proceden mayoritariamente de escenas de
        otras regiones del mundo. La brecha de dominio documentada en
        \cite{Lin_2025_CVPR} es geográfica y de escena; no hay evidencia
        de que sea neutral respecto de la población retratada.
\end{itemize}

\cajaaviso{Por qué el sesgo importa aquí, y por qué importa menos que en otros sistemas}{
Conviene ser preciso para no exagerar el riesgo ni minimizarlo. El sistema
no toma decisiones sobre personas individuales: no deniega embarques ni
señala sospechosos. Un error de conteo no perjudica directamente a la
persona mal contada. Sin embargo, si el subconteo se concentra en un
grupo, las zonas del terminal donde ese grupo se concentra reciben
sistemáticamente menos recursos --- menos personal, menos apertura de
mostradores. El daño no es individual sino distributivo, y es real aunque
sea indirecto.}

De aquí se sigue una recomendación concreta al frente técnico: que la
campaña de validación en sitio incluya una evaluación de desempeño
desagregada, al menos por franja de estatura y por presencia de equipaje
voluminoso. La recomendación es condicional en un sentido estricto ---
sin acceso a material del terminal no es realizable ---, y esa
dependencia refuerza la solicitud de acceso planteada en la
Sección~\ref{sec:auditoria}.

\subsubsection{Efecto sobre el comportamiento de las personas observadas}

Existe un impacto ético que no depende de que el sistema funcione bien y
que conviene dejar planteado con la cautela que corresponde. Se sostiene
habitualmente que saberse observado modifica la conducta; este informe no
puede afirmarlo como hecho acreditado, porque las únicas referencias
validadas del proyecto son los tres artículos técnicos de antecedentes y
no se ha incorporado literatura sobre efectos conductuales de la
videovigilancia. Formulado en condicional: si ese efecto existe en la
magnitud que la discusión pública le atribuye, en un aeropuerto ---
donde la vigilancia ya está normalizada --- el efecto marginal de añadir
un sistema de conteo sería previsiblemente pequeño, aunque no nulo, y la
normalización misma formaría parte del fenómeno. Sostener el punto con
firmeza exigiría respaldo bibliográfico que este informe todavía no
tiene.

% =====================================================================
\subsection{Impacto social}
\label{sec:impacto-social}

\subsubsection{Beneficios esperados}

El beneficio social principal es la seguridad. Las aglomeraciones densas
en espacios cerrados son un factor de riesgo conocido, y un sistema que
las detecta antes de que alcancen niveles críticos habilita una respuesta
preventiva en lugar de reactiva. En un terminal aéreo, donde la
concentración puede formarse rápidamente ante una demora o una
cancelación, el margen temporal que ofrece la detección temprana es el
beneficio más tangible.

Un beneficio secundario es la equidad en la asignación de recursos: la
medición sistemática permite fundamentar decisiones de dotación de
personal en evidencia y no en percepción, lo que tiende a favorecer a las
zonas y horarios con menor visibilidad.

Ambos beneficios se afirman aquí pero no se dimensionan, y así deben
presentarse al cliente. Su magnitud depende de la situación de partida en
el terminal --- si existen incidentes documentados por aglomeración y
cómo se decide hoy la dotación de personal ---, y esa es información que
solo LAP puede aportar. Mientras no la aporte, cualquier cifra de
beneficio sería una estimación sin base.

\subsubsection{Distribución de beneficios y cargas}

Una observación que conviene registrar por honestidad analítica: los
beneficios y las cargas de este sistema no recaen sobre los mismos
sujetos en la misma proporción.

\begin{center}
\small
\begin{tabularx}{\linewidth}{@{}L{3.2cm}YY@{}}
\toprule
\textbf{Actor} & \textbf{Beneficio} & \textbf{Carga} \\
\midrule
Pasajero &
Indirecto: menor tiempo de espera, mayor seguridad. &
Directa: su imagen es captada; no eligió participar. \\
\addlinespace
LAP (operador) &
Directo: eficiencia operativa, evidencia para decisiones, cumplimiento
de estándares de servicio. &
Indirecta: costo del sistema y responsabilidad por el tratamiento. \\
\addlinespace
Personal del terminal &
Apoyo a la toma de decisiones operativas. &
Riesgo de que la métrica se use para evaluación de desempeño individual,
finalidad distinta de la declarada. \\
\bottomrule
\end{tabularx}
\end{center}

La carga del pasajero es directa y el beneficio es indirecto; para el
operador la relación se invierte. Esto no invalida el proyecto --- es
la estructura habitual de las medidas de seguridad colectiva --- pero es
exactamente la razón por la cual el examen de proporcionalidad y el deber
de información no son formalidades.

La última fila merece atención específica: si los conteos por zona se
cruzan con los turnos del personal, el sistema se convierte de facto en
una herramienta de evaluación de desempeño laboral, con implicancias
laborales propias.

\cajaaviso{Riesgo a validar con el cliente}{El uso de los datos del
sistema para evaluar el desempeño del personal del terminal constituiría
una finalidad distinta de la declarada, con su propia base de
legitimación y su propio deber de información hacia los trabajadores.
Debe excluirse expresamente o tratarse por separado. Este frente propone
incorporar esa exclusión a la declaración de finalidades excluidas de la
Sección~\ref{sec:impacto-etico}, junto con las relativas a los
pasajeros.}

\subsubsection{Riesgo de dependencia y de exceso de confianza}

Un riesgo social poco discutido es la sustitución del criterio humano.
Si el personal de operaciones llega a confiar en el tablero más que en su
propia observación, un fallo del sistema en un momento crítico produce un
efecto peor que su ausencia, porque la atención humana ya se retiró.

La medida M-10 de la Sección~\ref{sec:privacidad-diseno} --- el sistema
alerta pero no decide --- responde a este riesgo, pero es insuficiente por
sí sola: requiere que los operadores conozcan las limitaciones del modelo,
incluido el modo de fallo por sobreestimación severa documentado en
\cite{Lin_2025_CVPR}. La capacitación del personal es una medida de
mitigación, no un accesorio.

% =====================================================================
\subsection{Impacto económico}
\label{sec:impacto-economico}

\cajaaviso{Nota sobre la ausencia de cifras}{Esta subsección identifica
categorías de costo y de beneficio, pero no las cuantifica. Hacerlo
requeriría información de LAP sobre infraestructura existente, costos de
personal y volumen de pasajeros que el equipo no posee. Una estimación
numérica sin esa base sería inventada, y este informe prefiere una
categoría sin cifra antes que una cifra sin sustento.}

\subsubsection{Costos}

\begin{center}
\small
\begin{tabularx}{\linewidth}{@{}L{3.3cm}YL{2.4cm}@{}}
\toprule
\textbf{Categoría} & \textbf{Contenido} & \textbf{Estado} \\
\midrule
Infraestructura de captura &
Cámaras adicionales o recalibración de las existentes. Depende
críticamente de si la infraestructura de CCTV actual es reutilizable. &
No cuantificado \\
\addlinespace
Cómputo de inferencia &
Servidores o dispositivos de borde. La medida M-01 (procesamiento local)
incrementa este rubro respecto de una solución en nube, y ese sobrecosto
es el precio del control de datos. &
No cuantificado \\
\addlinespace
Anotación de datos &
Etiquetado de material del terminal para adaptación de dominio. Es el
rubro que el enfoque de \cite{Lin_2025_CVPR} reduce de forma
sustancial. &
No cuantificado \\
\addlinespace
Cumplimiento &
Elaboración de documentación, señalización, auditorías, eventual
designación de ODP. &
No cuantificado \\
\addlinespace
Operación y mantenimiento &
Monitoreo del modelo, reentrenamiento periódico, soporte. &
No cuantificado \\
\bottomrule
\end{tabularx}
\end{center}

De este cuadro se desprende cuál es la pregunta económica decisiva, y
conviene formularla explícitamente porque condiciona todas las demás: si
la infraestructura de CCTV existente en el terminal es reutilizable, el
costo del proyecto se concentra en cómputo y desarrollo; si no lo es,
cambia de orden de magnitud. Es una pregunta que solo LAP puede
responder, y hasta que lo haga ninguna de las categorías anteriores
admite una cifra.

\subsubsection{Beneficios económicos}

Los beneficios plausibles son la optimización de la dotación de personal
mediante asignación basada en demanda medida; la reducción de tiempos de
espera, con su efecto sobre indicadores de calidad de servicio y sobre el
gasto de los pasajeros en zonas comerciales; y el valor de evitar
incidentes, que es el rubro de mayor magnitud potencial y el más difícil
de estimar, por tratarse de eventos de baja probabilidad y alto impacto.

\subsubsection{Costo del incumplimiento}

Un análisis económico que omitiera el costo del incumplimiento estaría
sesgado. Debe considerarse la exposición a sanciones administrativas de
la autoridad de protección de datos, el costo reputacional de un
incidente de privacidad en una infraestructura de alta visibilidad
pública, y el costo de rehacer el sistema si se detecta una no
conformidad después del despliegue.

Este informe no consigna los rangos de sanción aplicables. El régimen
sancionador de la Ley 29733 y de su reglamento se expresa en UIT, y
reproducir un rango o un importe sin haberlo contrastado contra el texto
vigente y contra el valor de la UIT del ejercicio correspondiente sería
introducir en el análisis económico exactamente el tipo de cifra sin
respaldo que el resto del documento evita. La exposición a sanción se
afirma, por tanto, como categoría de costo y no como monto.

Este último punto tiene una implicancia de secuenciación que conviene
subrayar: el costo de incorporar los requisitos de privacidad en la fase
de diseño es una fracción del costo de reconstruir un sistema desplegado.
La privacidad desde el diseño es, además de una obligación, la opción
económicamente racional.

% =====================================================================
\subsection{Impacto ambiental}
\label{sec:impacto-ambiental}

\cajaaviso{Nota sobre la ausencia de mediciones}{El equipo no ha medido
el consumo energético del entrenamiento ni de la inferencia. Las
consideraciones que siguen son cualitativas y estructurales. Se indica
expresamente qué mediciones permitirían convertirlas en cuantitativas.}

\subsubsection{Huella del entrenamiento}

El entrenamiento de modelos de visión es intensivo en cómputo y, por
tanto, en energía. Tres decisiones de diseño ya adoptadas o previstas
tienen efecto directo sobre esta huella:

\begin{itemize}[leftmargin=1.5em,itemsep=0.3em]
  \item \textbf{Reutilización de modelos preentrenados.} CSRNet parte de
        capas de VGG-16 ya entrenadas \cite{Li_2018_CVPR}, evitando
        entrenar desde cero.
  \item \textbf{Eficiencia de etiquetado.} El enfoque semi-supervisado de
        \cite{Lin_2025_CVPR} reduce el trabajo de anotación, aunque
        conviene no sobrevender el punto: reduce el esfuerzo humano, no
        necesariamente el cómputo, pues el esquema maestro-estudiante
        implica dos modelos en entrenamiento.
  \item \textbf{Eliminación del algoritmo húngaro.} Los autores de
        \cite{Lin_2025_CVPR} señalan expresamente que P2R prescinde de
        ese componente, que identifican como costoso en tiempo dentro de
        la estrategia de emparejamiento punto a punto.
\end{itemize}

Convertir esta subsección de cualitativa en cuantitativa depende de una
sola medición: las horas de GPU efectivamente empleadas durante el
desarrollo. Es un registro sencillo, pero tiene una propiedad que obliga
a pedirlo cuanto antes al frente técnico: no es reconstruible a
posteriori. Si no se lleva desde el inicio del entrenamiento, la huella
del desarrollo queda definitivamente sin cuantificar.

\subsubsection{Huella de la inferencia continua}

A escala de la vida útil del sistema, el consumo de la inferencia
continua supera previsiblemente al del entrenamiento: el entrenamiento
ocurre unas pocas veces, la inferencia se ejecuta permanentemente sobre
múltiples flujos de video.

Dos decisiones de diseño con motivación distinta convergen en reducir
esta huella:

\begin{itemize}[leftmargin=1.5em,itemsep=0.3em]
  \item \textbf{Frecuencia de muestreo.} Las aglomeraciones se forman en
        escalas de minutos. Procesar un fotograma cada varios segundos en
        lugar de a velocidad completa reduce el cómputo en un factor
        elevado sin pérdida operativa relevante. La misma decisión reduce
        el volumen de imágenes tratadas, por lo que sirve simultáneamente
        al principio de minimización.
  \item \textbf{Procesamiento en el borde.} La medida M-01, adoptada por
        razones de protección de datos, evita transmitir video de forma
        continua hacia la nube, con el ahorro de transmisión asociado.
\end{itemize}

\cajaaviso{Convergencia entre minimización y sostenibilidad}{Es un
resultado analítico útil de esta sección: las decisiones que reducen la
exposición de datos personales --- muestrear menos, procesar localmente,
no persistir imágenes, agregar la salida --- son las mismas que reducen
el consumo de cómputo, de transmisión y de almacenamiento. Ambos
principios apuntan en la misma dirección, y conviene presentarlos así
ante el cliente porque refuerza el argumento de privacidad con una
justificación adicional que no depende del cumplimiento normativo.}

\subsubsection{Hardware y residuos electrónicos}

Si el proyecto requiere cámaras o servidores nuevos, su fabricación y su
eventual disposición final integran la huella. Reutilizar la
infraestructura de CCTV existente, además de su efecto económico
(Sec.~\ref{sec:impacto-economico}), es la opción ambientalmente
preferible.

Queda un punto por confirmar con el cliente antes de cerrar esta
subsección: si LAP cuenta con una política institucional de
sostenibilidad o con compromisos de reporte ambiental. De existir, este
apartado debería alinearse con ellos y utilizar sus propias categorías de
medición, en lugar de proponer un marco paralelo.

% =====================================================================
\subsection{Síntesis del análisis de impacto}
\label{sec:sintesis-impacto}

Integrando las cuatro dimensiones, el perfil de impacto del sistema puede
resumirse así.

{\footnotesize
\begin{longtable}{@{}L{1.9cm}L{4.0cm}L{4.0cm}L{2.6cm}@{}}
\toprule
\textbf{Dimensión} & \textbf{Impacto positivo} & \textbf{Riesgo principal} & \textbf{Mitigación} \\
\midrule
\endfirsthead
\toprule
\textbf{Dimensión} & \textbf{Impacto positivo} & \textbf{Riesgo principal} & \textbf{Mitigación} \\
\midrule
\endhead
\bottomrule
\endfoot
Ética &
Alternativa menos invasiva que la identificación biométrica para la misma
finalidad. &
Deslizamiento de finalidad hacia identificación o seguimiento. &
Declaración de finalidades excluidas; M-02, M-03. \\
\addlinespace
Social &
Prevención de incidentes por aglomeración; asignación de recursos basada
en evidencia. &
Distribución desigual de beneficio y carga; uso laboral no declarado;
exceso de confianza en el sistema. &
Deber de información; exclusión expresa del uso laboral; capacitación;
M-10. \\
\addlinespace
Económica &
Eficiencia operativa; valor de incidentes evitados. &
Costo de reconstrucción si se detecta no conformidad tras el
despliegue. &
Privacidad desde el diseño; auditoría previa
(Sec.~\ref{sec:auditoria}). \\
\addlinespace
Ambiental &
Menor huella que alternativas de mayor granularidad; reutilización de
modelos preentrenados. &
Consumo de la inferencia continua a lo largo de la vida útil. &
Frecuencia de muestreo reducida; procesamiento en el borde;
reutilización de hardware. \\
\end{longtable}}

Tres conclusiones transversales se desprenden del análisis.

\begin{enumerate}[leftmargin=1.6em,itemsep=0.4em]
  \item \textbf{Las decisiones de minimización operan simultáneamente en
        las cuatro dimensiones.} No persistir imágenes, agregar la salida
        y muestrear a baja frecuencia reducen a la vez el riesgo ético,
        la exposición social, el costo de almacenamiento y la huella
        ambiental. Son el núcleo de la propuesta de este frente
        precisamente porque ningún objetivo compite con otro.
  \item \textbf{El riesgo dominante es de gobernanza, no de tecnología.}
        El deslizamiento de finalidad y el uso no declarado no se
        previenen con mejor ingeniería sino con documentación,
        compromisos contractuales y auditoría. Esto significa que el
        entregable de este frente --- la documentación de cumplimiento
        --- no es accesorio al proyecto técnico: es la mitigación
        principal del riesgo principal.
  \item \textbf{Varios impactos permanecen sin cuantificar por falta de
        información del cliente y de mediciones propias.} Se ha optado
        por declararlo antes que por estimar sin base. La lista completa
        de lo pendiente, con la indicación de quién puede cerrar cada
        punto, está en \texttt{PENDIENTES.md}, y el estado de cumplimiento
        en \texttt{CHECKLIST\_\allowbreak CUMPLIMIENTO.md}.
\end{enumerate}

La segunda de esas conclusiones merece una nota final, porque tiene una
consecuencia sobre cómo se presenta el proyecto. Si el riesgo dominante
es de gobernanza, entonces el trabajo de documentación de cumplimiento no
es un requisito administrativo del curso sino la mitigación principal del
riesgo principal, y así debería exponerse ante el cliente. Este frente
sostiene ese argumento y recomienda presentarlo de forma explícita.
